% ==========
%  CHƯƠNG 9
% ==========
\OpenGrid

\ParallelChapter
{Equity}
{Công bằng}
{}

\TriPara
{
\EN{
\begin{flushright}
  \emph{Mathematical reasoning and sense making must be evident in the mathematical experiences of all students.}
\end{flushright}}
}
{
\VI{
\begin{flushright}
  \emph{Luận toán học và hiểu phải hiện diện rõ trong trải nghiệm toán học của mọi học sinh.}
\end{flushright}}
}
{
  \NOTE{Xác lập lập trường xuyên suốt của chương: luận toán học và hiểu không phải là đặc quyền của nhóm học sinh có năng lực cao, cũng không phải là phần bổ sung sau khi đã học xong nội dung. Chúng phải hiện diện trong chính trải nghiệm toán học của mọi học sinh.}
}

\TriPara
  {
    \EN{In this publication, we build on the concept of equity as defined in \emph{Principles and Standards for School Mathematics} (NCTM 2000a), in which equity is described in terms of having high expectations of, and providing support for, all students. This population includes students of all races and cultures; students with low socioeconomic backgrounds; students whose first language is not English; both girls and boys; both gifted and low-achieving students; and students with learning disabilities, emotional problems, and behavioral problems. All too frequently, one or another of these groups is deemed too difficult to teach. Often, well-meaning teachers and administrators cause inequity through practices that create biases they do not intend. High schools can monitor equity by paying focused attention to phenomena that pose potential barriers to engaging every student in reasoning and sense making. These phenomena include the following:\vspace{4\baselineskip}

    \begin{itemize}
      \item The \emph{courses} that students take have an impact on the opportunities that they have for reasoning and sense making.
      
      \item \emph{Students' demographics} too often predict the opportunities students have for reasoning and sense making.
      
      \item \emph{Expectations, beliefs, and biases} have an impact on the mathematical learning opportunities provided for students.
    \end{itemize}
    }
  }
  {
    \VI{Trong ấn phẩm này, quan niệm về công bằng được triển khai trên cơ sở cách hiểu đã được xác định trong \emph{Principles and Standards for School Mathematics} (NCTM 2000a), theo đó công bằng được mô tả qua việc đặt kì vọng cao đối với mọi học sinh và đồng thời cung cấp sự hỗ trợ cần thiết cho tất cả các em. Nhóm học sinh ấy bao gồm học sinh thuộc mọi chủng tộc và nền văn hóa; học sinh có hoàn cảnh kinh tế - xã hội khó khăn; học sinh mà tiếng Anh không phải là ngôn ngữ thứ nhất; cả nữ sinh lẫn nam sinh; cả học sinh có năng khiếu lẫn học sinh có thành tích thấp; cùng những học sinh gặp khó khăn trong học tập, có vấn đề về cảm xúc hoặc có vấn đề về hành vi. Quá thường xuyên, nhóm này hay nhóm khác trong số đó bị xem là quá khó để dạy. Không ít khi, chính những giáo viên và nhà quản lí vốn có thiện chí lại tạo ra sự bất công thông qua những thực hành làm nảy sinh các thiên lệch mà bản thân họ không hề chủ ý. Các trường trung học có thể theo dõi mức độ công bằng bằng cách tập trung chú ý đến những hiện tượng có thể trở thành rào cản đối với việc lôi cuốn mọi học sinh tham gia vào luận và hiểu. Những hiện tượng ấy gồm có:

    \begin{itemize}
      \item \emph{Khóa học} mà học sinh theo học tác động đến cơ hội mà các em có được để phát triển luận và hiểu.
      
      \item \emph{Đặc điểm nhân khẩu học} của học sinh quá thường xuyên dự báo trước cơ hội mà các em có được để phát triển luận và hiểu.
      
      \item \emph{Kì vọng, niềm tin và thiên lệch} tác động đến cơ hội học toán được tạo ra cho học sinh.
    \end{itemize}}
  }
  {
    \NOTE{Đoạn này xác định công bằng không chỉ là quyền tiếp cận chung chung, mà là quyền được tham gia vào những trải nghiệm học toán có luận và hiểu. Ba hiện tượng được liệt kê ở cuối đoạn cũng là ba trục phân tích chính của chương: khóa học, đặc điểm nhân khẩu học, và kì vọng - niềm tin - thiên lệch trong nhà trường.}
  }

\ParallelSection
  {Courses}
  {Khóa học}
  {}

\TriPara
  {
    \EN{Over the past two decades, enrollment in advanced and college-preparatory mathematics courses at the high school level has increased noticeably (Planty, Provasnik, and Daniel 2007). At the same time, high schools have continued to employ some form of tracking or ability grouping (Hallinan 2004). With the tracking of students, high schools run the risk of promoting inequitable opportunities for mathematical learning for those students placed in remedial classes or in classes on “low” or “regular” tracks. For example, teacher quality varies across tracks, with more out-of-field teachers assigned to teach low-track classes (Tate and Rousseau 2002). High schools that use tracking or ability grouping can be proactive by diligently ensuring that at every level, all mathematics courses offer students opportunities for reasoning and sense making. In addition, with proper school support and access to such opportunities, students do not have to be mired in one track for their entire high school career. Successful high schools make sure that moves to different tracks are both offered and actively supported (Education Trust 2005). Moreover, we find numerous success stories of schools that have eliminated tracking with very positive results, including gains in mathematics course taking and achievement (Garrity 2004) and large numbers of minority students going on to attend college (Alvarez and Mehan 2006). An important point to note is that these schools' success was not attributed solely to eliminating tracking but to other elements of support for students and teachers in the schools, as well.}
  }
  {
    \VI{Trong hai thập niên qua, việc ghi danh vào khóa toán nâng cao và khóa toán dự bị đại học ở bậc trung học phổ thông đã tăng lên rõ rệt (Planty, Provasnik, and Daniel 2007). Đồng thời, các trường trung học vẫn tiếp tục áp dụng hình thức phân luồng hoặc phân nhóm theo năng lực (Hallinan 2004). Khi phân luồng học sinh như vậy, các trường trung học có nguy cơ tạo ra cơ hội học toán không công bằng đối với học sinh được xếp vào lớp phụ đạo hoặc vào lớp thuộc luồng ``thấp'' hoặc luồng ``thường''. Chẳng hạn, chất lượng giáo viên khác nhau giữa các luồng, trong đó lớp thuộc luồng thấp thường có nhiều giáo viên dạy trái chuyên môn hơn (Tate and Rousseau 2002). Các trường trung học sử dụng hình thức phân ban hoặc phân nhóm theo năng lực có thể chủ động ứng phó bằng cách bảo đảm một cách nghiêm túc rằng ở mọi trình độ, mọi khóa toán đều đem lại cho học sinh cơ hội luận và hiểu. Ngoài ra, nếu có sự hỗ trợ thích đáng từ phía nhà trường và có điều kiện tiếp cận cơ hội như vậy, học sinh không nhất thiết phải bị kẹt mãi trong một luồng suốt những năm học trung học phổ thông. Các trường trung học thành công bảo đảm rằng việc chuyển sang luồng khác không chỉ được mở ra như khả năng mà còn được hỗ trợ một cách tích cực (Education Trust 2005). Hơn nữa, nhiều câu chuyện thành công cho thấy có những trường đã xóa bỏ phân luồng và đạt được kết quả rất tích cực, bao gồm việc học sinh theo học nhiều khóa toán hơn và đạt thành tích cao hơn (Garrity 2004), cũng như việc có số lượng lớn học sinh thuộc nhóm thiểu số tiếp tục theo học đại học (Alvarez and Mehan 2006). Điểm quan trọng cần lưu ý là thành công của các trường này không chỉ được quy cho riêng việc xóa bỏ phân luồng, mà còn gắn với những yếu tố hỗ trợ khác dành cho học sinh và giáo viên trong nhà trường.}
  }
  {
    \NOTE{Nên giữ nhất quán cặp thuật ngữ \emph{tracking/track} là ``phân luồng/luồng'' trong cụm này. Trọng tâm của đoạn không phải là phản đối riêng việc phân nhóm học sinh, mà là cảnh báo nguy cơ bất công khi phân luồng làm cho học sinh ở luồng thấp có ít cơ hội tham gia vào luận và hiểu, ít giáo viên đúng chuyên môn hơn, và khó chuyển sang luồng học tập khác nếu thiếu hỗ trợ chủ động từ nhà trường.}
  }

\TriPara
  {
    \EN{Schools with a curricular structure in which multiple levels of the same class are offered (e.g., Informal Geometry vs. Honors Geometry or Honors Integrated Mathematics vs. Integrated Mathematics) have a responsibility to engage all students in every class in reasoning and sense making, thereby ensuring that the students are learning the same major concepts regardless of the course title.}
  }
  {
    \VI{Các trường học có cấu trúc chương trình trong đó cùng một khóa học được tổ chức ở nhiều cấp độ khác nhau (ví dụ: hình học không hình thức so với hình học nâng cao, hoặc toán tích hợp nâng cao so với toán tích hợp) có trách nhiệm bảo đảm rằng mọi học sinh trong mọi lớp học đều được tham gia vào luận và hiểu. Qua đó, nhà trường bảo đảm rằng học sinh đều đang học cùng những khái niệm lớn cốt lõi, bất kể tên gọi của khóa học là gì.}
  }{\NOTE{}}

\TriPara
  {
    \EN{Students' unpreparedness for algebra is a current problem for high schools and often leads to schools' needing to provide some sort of remediation or additional academic support. Schools deal with this problem in various ways, such as stretching algebra over two years or placing students in a prealgebra class. These strategies can be counterproductive, especially if the result is that students are confined to the low track with limited attention to reasoning for the duration of high school and without the opportunity to catch up to their peers. High schools can support students who are underprepared for first-year algebra without limiting their future progress. For example, some high schools take early action with struggling students by offering courses in the summer before high school or as part of after-school programs (Education Trust 2005). Other schools provide first-year algebra students with general academic support in the form of teacher and peer tutoring, Saturday programs, and double course periods for some courses (Huebner, Corbett, and Phillippo 2006). In all instances, these high schools provide support early on for those students entering high school below grade level. Although programs to support underprepared students require additional resources, schools must make additional instruction to assist these students a priority. Algebra presents great opportunities for sense making and reasoning and lays a foundation for further study in mathematics. Thus, taking steps to ensure that all first-year algebra students are supported in ways that allow them to become proficient with algebraic concepts is of crucial importance.}
  }
  {
    \VI{Tình trạng học sinh chưa sẵn sàng cho đại số hiện là vấn đề ở các trường trung học phổ thông và thường khiến nhà trường phải tổ chức một số hình thức phụ đạo hoặc hỗ trợ học tập bổ sung. Các trường xử lí vấn đề này theo nhiều cách khác nhau, chẳng hạn như kéo dài chương trình đại số thành hai năm hoặc xếp học sinh vào lớp tiền đại số. Tuy nhiên, những chiến lược ấy có thể phản tác dụng, nhất là khi kết quả là học sinh bị giữ lại ở luồng thấp trong suốt thời gian học trung học phổ thông, ít được chú trọng đến luận, và không có cơ hội bắt kịp các bạn cùng trang lứa. Các trường trung học phổ thông có thể hỗ trợ học sinh chưa sẵn sàng cho đại số năm đầu mà không làm hạn chế sự tiến bộ sau này của các em. Chẳng hạn, một số trường can thiệp sớm đối với học sinh gặp khó khăn bằng cách mở khóa học vào mùa hè trước khi vào trung học hoặc tổ chức chương trình sau giờ học (Education Trust 2005). Những trường khác hỗ trợ học sinh học đại số năm đầu bằng hình thức hỗ trợ học tập chung, như được giáo viên và bạn học kèm cặp, chương trình học vào ngày thứ Bảy, hoặc tăng số tiết cho một số môn học (Huebner, Corbett, and Phillippo 2006). Trong tất cả trường hợp ấy, các trường này đều cung cấp sự hỗ trợ từ sớm cho học sinh bước vào trung học khi chưa đạt mức yêu cầu tương ứng. Mặc dù chương trình hỗ trợ học sinh chưa sẵn sàng đòi hỏi thêm nguồn lực, nhà trường vẫn phải coi việc tăng cường giảng dạy để hỗ trợ những học sinh này là ưu tiên. Đại số mở ra nhiều cơ hội cho hiểu và luận, đồng thời đặt nền móng cho việc học toán về sau. Vì vậy, việc thực hiện các biện pháp nhằm bảo đảm rằng tất cả học sinh học đại số năm đầu đều được hỗ trợ theo những cách giúp các em nắm vững các khái niệm đại số là điều hết sức quan trọng.}
  }{\NOTE{}}

\TriPara
  {
    \EN{The increasing emphasis on assessment in schools under the No Child Left Behind (NCLB) Act of 2001 (Public Law 107--110) has important consequences for opportunities for students to engage in reasoning and sense making in the classroom. In most instances, district and school policies on assessment heavily influence decisions about curriculum in the classroom. With extreme pressure for students to perform well on state assessments, teachers find themselves spending more time teaching test-taking skills and reviewing or covering those topics that are likely to be on the test. This emphasis, coupled with the fact that often high-stakes tests focus on students' ability to do procedural tasks rather assess their reasoning and sense making, equates to students' spending less time in class on developing their reasoning and sense-making abilities. This effect is even more prevalent for teachers in low-performing, high-minority schools who, under NCLB, often have even more pressure to ensure that their students succeed (Tate and Rousseau 2006). Coherence among assessment, curricula, and instruction is instrumental to addressing this issue. A discussion about the importance of coherence is included in chapter 11.}
  }
  {
    \VI{Việc ngày càng nhấn mạnh đến đánh giá trong nhà trường dưới tác động của Đạo luật No Child Left Behind (NCLB) năm 2001 (Public Law 107--110) kéo theo hệ quả quan trọng đối với cơ hội để học sinh tham gia vào luận và hiểu trong lớp học. Trong phần lớn trường hợp, chính sách đánh giá ở cấp học khu và nhà trường có ảnh hưởng mạnh mẽ đến quyết định về chương trình học trong lớp. Trước áp lực rất lớn buộc học sinh phải đạt kết quả tốt trong kì đánh giá cấp bang, giáo viên thường phải dành nhiều thời gian hơn cho việc dạy kĩ năng làm bài kiểm tra, cũng như ôn tập hoặc dạy những chủ đề có khả năng xuất hiện trong bài kiểm tra. Sự nhấn mạnh này, kết hợp với thực tế là bài kiểm tra mang tính quyết định cao thường tập trung vào khả năng thực hiện nhiệm vụ mang tính thủ tục hơn là đánh giá luận và hiểu của học sinh, khiến học sinh có ít thời gian hơn trên lớp để phát triển năng lực luận và hiểu. Tác động này còn phổ biến hơn đối với giáo viên ở các trường có kết quả thấp và tỉ lệ học sinh thiểu số cao; dưới khuôn khổ của NCLB, họ thường phải chịu áp lực lớn hơn nữa để bảo đảm học sinh của mình thành công (Tate and Rousseau 2006). Sự nhất quán giữa đánh giá, chương trình học và giảng dạy giữ vai trò then chốt trong việc giải quyết vấn đề này. Một phần bàn luận về tầm quan trọng của sự nhất quán ấy được trình bày trong Chương 11.}
  }{\NOTE{}}

\TriPara
  {
    \EN{To ensure that this publication's vision becomes a reality, teachers must hold high expectations for all students and find strategies to ensure that all students, at any level and in any course, can reason mathematically by engaging in challenging mathematical tasks. One such strategy is using problems that have multiple entry points so that students at different levels of mathematical experience and with different interests can all engage meaningfully in reasoning about the problem. Using problems with multiple entry points allows students to use their individual mathematical strengths in approaching the problem and gives them opportunities to make sense of the reasoning of other students.}
  }
  {
    \VI{Để bảo đảm rằng tầm nhìn của ấn phẩm này trở thành hiện thực, giáo viên phải đặt kì vọng cao đối với mọi học sinh và tìm ra chiến lược để bảo đảm rằng tất cả học sinh, ở mọi trình độ và trong bất kì khóa học nào, đều có thể luận toán học thông qua việc tham gia các nhiệm vụ toán học giàu thách thức. Một chiến lược như vậy là dùng bài toán có nhiều điểm vào, để học sinh với mức độ trải nghiệm toán học và mối quan tâm khác nhau đều có thể tham gia một cách có ý nghĩa vào việc luận về bài toán. Việc dùng bài toán có nhiều điểm vào cho phép học sinh vận dụng điểm mạnh toán học riêng của mình khi tiếp cận bài toán, đồng thời tạo cơ hội cho các em hiểu cách luận của học sinh khác.}
  }{\NOTE{}}

\TriPara
  {
    \EN{Several examples in section 2 of this publication describe tasks that have multiple entry points. For instance, example 7, ``Finding Balance,'' presents a problem that can be solved with or without the use of equations. If students are able to use equations flexibly, as student 1 can in the example, they can represent the described situation symbolically, solve the equation, and find the solution. However, other students who are either less confident, have had limited experience with equations, or are unable to use them flexibly might approach the problem by first drawing a picture to represent the scenario, then move toward writing an equation. Others might use a more informal reasoning approach, such as the one used by student 2, to find the answer. Exploring multiple approaches to solving a problem can forge important connections among different mathematical domains and strengthen the reasoning abilities of all learners. In example 7, student 1's solution may help student 2 consider a more formal approach, and student 2's solution may help student 1 make sense of formal equation solving. For this reason, students should be given opportunities to discuss their different approaches with one another and work toward understanding how each approach relates to the others.}
  }
  {
    \VI{Một số ví dụ trong Phần 2 của ấn phẩm này mô tả những nhiệm vụ có nhiều điểm vào. Chẳng hạn, Ví dụ 7, ``Tìm thế cân bằng'', đưa ra bài toán có thể giải bằng phương trình hoặc không cần dùng phương trình. Nếu học sinh có thể sử dụng phương trình một cách linh hoạt, như Học sinh 1 trong ví dụ, thì các em có thể biểu diễn tình huống được mô tả bằng kí hiệu, giải phương trình và tìm ra đáp án. Tuy nhiên, những học sinh khác---những em kém tự tin hơn, có ít kinh nghiệm với phương trình, hoặc chưa thể sử dụng phương trình linh hoạt---có thể tiếp cận bài toán bằng cách trước hết vẽ hình để biểu diễn tình huống, rồi từ đó chuyển dần sang việc viết phương trình. Những học sinh khác nữa có thể dùng một cách luận không chính thức hơn, như cách Học sinh 2 đã dùng, để tìm ra đáp án. Việc tìm hiểu nhiều cách khác nhau để giải bài toán có thể tạo ra mối liên hệ quan trọng giữa các lĩnh vực toán học khác nhau và củng cố năng lực luận của mọi người học. Trong Ví dụ 7, lời giải của Học sinh 1 có thể giúp Học sinh 2 nghĩ tới cách tiếp cận hình thức hơn, còn lời giải của Học sinh 2 có thể giúp Học sinh 1 hiểu được việc giải phương trình hình thức. Vì vậy, học sinh cần được tạo cơ hội để trao đổi với nhau về các cách tiếp cận khác nhau của mình và cùng đi tới chỗ hiểu các cách tiếp cận ấy liên hệ với nhau như thế nào.}
  }{\NOTE{}}

\TriPara
  {
    \EN{In any given high school classroom, students manifest varied levels of readiness for mathematical content, both among those who need extra support and among those who are ready to go further than the rest of the class. Often, students are deemed ``gifted'' or mathematically ``talented'' on the basis of a series of tests. However, teachers and schools should realize that students' mathematical ``talent'' is fluid (Samuels 2008) and that although a student might show great promise in geometry, for example, she or he might need extra support with algebraic content. For those students who demonstrate readiness to move beyond the core curriculum, teachers can provide opportunities to delve more deeply into the content being studied by the rest of the class. These students must have opportunities for reasoning and sense making that maximize their mathematical learning experiences. Although acceleration is definitely one component, opportunities for enrichment, advanced content, and more in-depth study can all be part of the mathematical education of these students, as well (National Association for Gifted Children 1994). For instance, in example 15, ``Circling the Points,'' students discuss the result that three noncollinear points in a plane determine a unique circle. As a next step for a student who is interested or is ready to move beyond the task posed to the rest of the class, the teacher might pose the questions ``Do four noncollinear points in the plane determine a unique circle?'' ``What are the conditions?'' ``What are some conjectures?'' This series of questions leads the student to thinking about inscribed angles, a topic that would typically not be discussed in the course until later. Thus, the student is asked to think more deeply about the original problem assigned to the class.}
  }
  {
    \VI{Trong bất kì lớp học trung học phổ thông nào, học sinh đều thể hiện mức độ sẵn sàng khác nhau đối với nội dung toán học, ở cả nhóm cần thêm hỗ trợ lẫn nhóm đã sẵn sàng đi xa hơn so với phần còn lại của lớp. Thông thường, học sinh được xem là ``có năng khiếu'' hoặc ``có tài năng toán học'' dựa trên kết quả của loạt bài kiểm tra. Tuy nhiên, giáo viên và nhà trường cần nhận thức rằng ``tài năng'' toán học của học sinh không cố định (Samuels 2008): học sinh có thể thể hiện tiềm năng nổi trội trong hình học, chẳng hạn, nhưng lại cần thêm hỗ trợ đối với nội dung đại số. Với những học sinh cho thấy đã sẵn sàng vượt ra ngoài chương trình cốt lõi, giáo viên có thể tạo cơ hội để các em đi sâu hơn vào chính nội dung mà cả lớp đang học. Những học sinh này cần có cơ hội luận và hiểu nhằm tối đa hóa trải nghiệm học toán của mình. Mặc dù học tăng tốc chắc chắn là một thành phần, nhưng cơ hội bồi dưỡng mở rộng, tiếp cận nội dung nâng cao và học chuyên sâu hơn cũng đều có thể là những phần cấu thành của giáo dục toán học dành cho các em (National Association for Gifted Children 1994). Chẳng hạn, trong Ví dụ 15, ``Tâm ở đâu?'', học sinh thảo luận kết quả rằng ba điểm không thẳng hàng trong mặt phẳng xác định duy nhất một đường tròn. Như một bước tiếp theo dành cho học sinh có hứng thú hoặc đã sẵn sàng đi xa hơn nhiệm vụ chung của lớp, giáo viên có thể đặt câu hỏi: ``Bốn điểm không thẳng hàng trong mặt phẳng có xác định đường tròn duy nhất không?'' ``Điều kiện là gì?'' ``Có thể nêu ra phỏng đoán nào?'' Chuỗi câu hỏi này dẫn học sinh suy nghĩ về góc nội tiếp, chủ đề thường chỉ được bàn tới ở giai đoạn sau của khóa học. Như vậy, học sinh được yêu cầu suy nghĩ sâu hơn về chính bài toán ban đầu được giao cho cả lớp.}
  }
  {
    \NOTE{Đoạn này mở rộng cách hiểu về công bằng: công bằng không chỉ là hỗ trợ học sinh gặp khó khăn, mà còn là tạo cơ hội thích đáng cho học sinh đã sẵn sàng đi xa hơn. Điểm then chốt là không xem ``tài năng'' toán học như phẩm chất cố định; học sinh có thể nổi trội ở một mảng nhưng vẫn cần hỗ trợ ở mảng khác, và ngược lại, có thể được đưa vào luận và hiểu sâu hơn thông qua nhiệm vụ mở rộng phù hợp.}
  }

\ParallelSection
  {Student Demographics and Opportunities for Learning}
  {Đặc điểm nhân khẩu học của học sinh và cơ hội học tập}
  {}

\TriPara
  {
    \EN{Discrepancies in achievement between student groups continue to raise concern for educators, families, and leaders in the United States. As reported on the mathematics portion of the 2003 National Assessment of Educational Progress (NAEP) (Lubienski and Crockett 2007), 52 percent of Hispanic and 61 percent of African American eighth-grade students scored below the basic level as compared with 20 percent of white eighth-grade students. Only 7 percent of African American students and 12 percent of Hispanic students scored at or above the proficient level in comparison with 37 percent of white students. An important finding of the same study was that white (91 percent) students were more likely to have fully credentialed teachers than African American (80 percent) and Hispanic (83 percent) students. This finding, coupled with the fact that students of teachers who majored in mathematics in college scored seven to fifteen points higher on NAEP than students of teachers who were not mathematics majors (Lubienski and Crockett 2007), indicates an important connection between discrepancies in resources and achievement.}
  }
  {
    \VI{Chênh lệch về thành tích giữa các nhóm học sinh tiếp tục gây quan ngại cho nhà giáo dục, gia đình và nhà lãnh đạo tại Hoa Kì. Theo báo cáo phần toán của Kì đánh giá quốc gia về tiến bộ giáo dục (National Assessment of Educational Progress - NAEP) năm 2003 (Lubienski and Crockett 2007), có tới 52 phần trăm học sinh lớp 8 người Hispanic và 61 phần trăm học sinh lớp 8 người Mĩ gốc Phi đạt điểm dưới mức cơ bản, so với 20 phần trăm học sinh da trắng lớp 8. Chỉ 7 phần trăm học sinh Mĩ gốc Phi và 12 phần trăm học sinh Hispanic đạt mức thành thạo trở lên, trong khi con số này ở học sinh da trắng là 37 phần trăm. Phát hiện quan trọng khác của cùng nghiên cứu cho thấy học sinh da trắng (91 phần trăm) có nhiều khả năng được học với giáo viên có đầy đủ chứng chỉ chuyên môn hơn so với học sinh Mĩ gốc Phi (80 phần trăm) và học sinh Hispanic (83 phần trăm). Phát hiện này, kết hợp với thực tế rằng học sinh của những giáo viên tốt nghiệp đại học chuyên ngành toán đạt điểm NAEP cao hơn từ bảy đến mười lăm điểm so với học sinh của giáo viên không chuyên toán (Lubienski and Crockett 2007), cho thấy mối liên hệ quan trọng giữa sự chênh lệch về nguồn lực và thành tích học tập.}
  }
  {
    \NOTE{Nên giữ \emph{Hispanic} trong bản dịch vì đây là nhóm nhân khẩu học theo cách phân loại ở Hoa Kì; dịch thành ``gốc Tây Ban Nha'' dễ làm hẹp nghĩa sang nguồn gốc từ Tây Ban Nha. Trọng tâm của đoạn là chỉ ra chênh lệch thành tích gắn với chênh lệch nguồn lực, đặc biệt là khả năng tiếp cận giáo viên có đủ chứng chỉ và có nền tảng chuyên môn toán.}
  }

\TriPara
  {
    \EN{Educators in the United States continue to grapple with economic inequity, which has a significant impact on districts' and schools' abilities to improve the opportunities to learn for the country's most underserved students (Tate and Rousseau 2007). Astounding discrepancies remain between the amount of money spent by the wealthiest districts as compared with that spent by the poorest, often high-minority, urban districts (Darling-Hammond 2004). These discrepancies exist between states as well as within states. Districts that spend more money often have smaller class sizes and more instructional resources, are able to offer a wider range of classes, and have teachers who are better paid and more experienced (Darling-Hammond 2004)}
  }
  {
    \VI{Các nhà giáo dục tại Hoa Kì vẫn đang phải vật lộn với bất công về kinh tế, yếu tố có tác động đáng kể đến khả năng của các học khu và nhà trường trong việc cải thiện cơ hội học tập cho học sinh bị thiệt thòi nhất của quốc gia (Tate and Rousseau 2007). Chênh lệch đáng kinh ngạc vẫn tồn tại giữa mức chi tiêu của các học khu giàu có nhất và mức chi tiêu của các học khu nghèo nhất, thường là học khu đô thị có tỉ lệ học sinh thiểu số cao (Darling-Hammond 2004). Chênh lệch này tồn tại không chỉ giữa các bang mà còn ngay trong từng bang. Học khu chi tiêu nhiều hơn thường có sĩ số lớp học nhỏ hơn và nhiều nguồn lực giảng dạy hơn, có khả năng cung cấp nhiều khóa học đa dạng hơn, đồng thời có đội ngũ giáo viên được trả lương cao hơn và giàu kinh nghiệm hơn (Darling-Hammond 2004).}
  }
  {
    \NOTE{Đoạn này chuyển trọng tâm từ chênh lệch thành tích sang chênh lệch nguồn lực kinh tế. Điểm chính không chỉ là học khu giàu chi tiêu nhiều hơn, mà là mức chi tiêu ấy gắn với sĩ số lớp, nguồn lực dạy học, phạm vi khóa học được cung cấp, mức lương và kinh nghiệm của giáo viên; từ đó tác động trực tiếp đến cơ hội học toán của học sinh.}
  }

\TriPara
  {
    \EN{Providing students with more opportunities to learn mathematics has serious financial consequences. For example, creating Saturday or summer programs to provide additional support to those students not ready for algebra requires additional money. Unfortunately, such programs are needed most in schools or districts that often have fewer resources. Similarly, schools that are predominantly low-income, African American, or Hispanic often contend with high teacher-turnover rates. High turnover has the unfortunate consequence that those students who are most in need of the best possible teachers are very often being taught by the least experienced teachers (Bishop and Forgasz 2006). In addition, teaching out of field is more prevalent in high-minority and lowsocioeconomic schools than in other schools, with the result that poor, minority students in the United States are often being taught by teachers who are underqualified (Education Trust 2008).}
  }
  {
    \VI{Việc tạo thêm cơ hội học toán cho học sinh kéo theo hệ quả lớn về tài chính. Chẳng hạn, việc tổ chức chương trình học vào ngày thứ Bảy hoặc trong mùa hè nhằm hỗ trợ bổ sung cho học sinh chưa sẵn sàng học đại số đòi hỏi phải có thêm kinh phí. Đáng tiếc là chương trình như vậy lại cần thiết nhất ở trường hoặc học khu vốn thường có ít nguồn lực hơn. Tương tự, các trường chủ yếu phục vụ học sinh thu nhập thấp, học sinh Mĩ gốc Phi hoặc học sinh Hispanic thường phải đối mặt với tỉ lệ giáo viên thay đổi rất cao. Hệ quả đáng tiếc của tình trạng này là học sinh cần đến giáo viên tốt nhất lại rất thường xuyên được dạy bởi giáo viên ít kinh nghiệm nhất (Bishop and Forgasz 2006). Bên cạnh đó, tình trạng giáo viên dạy trái chuyên môn xảy ra phổ biến hơn ở trường có ti lệ học sinh thiểu số cao và điều kiện kinh tế - xã hội thấp so với trường khác, dẫn đến việc học sinh nghèo và học sinh thuộc nhóm thiểu số tại Hoa Kì thường được dạy bởi giáo viên chưa đủ chuẩn chuyên môn (Education Trust 2008).}
  }
  {
    \NOTE{Đoạn này làm rõ mặt tài chính của công bằng: mở thêm cơ hội học toán, chẳng hạn chương trình thứ Bảy hoặc mùa hè, đòi hỏi nguồn lực; nhưng chính những trường cần các chương trình này nhất lại thường thiếu nguồn lực nhất. Hệ quả là học sinh vốn cần giáo viên tốt nhất lại dễ gặp giáo viên ít kinh nghiệm hoặc dạy trái chuyên môn.}
  }

\TriPara
  {
    \EN{The vision set forth for high school mathematics in this publication depends on teachers who have strong knowledge of mathematical content. Without this knowledge and the confidence it inspires, teachers have difficulty pushing reasoning to the forefront of the curriculum. In essence, more money, resources, and better teachers are being provided to the students who come from the wealthiest families, whereas less money, fewer resources, and less-qualified teachers are provided for the students from the poorest families (Darling-Hammond 2004). Simply providing equitable resources for all students is certainly not sufficient to ensure that each student will develop reasoning and sense making in mathematics, but it is surely necessary for real progress to be made (Bishop and Forgasz 2006). Therefore, attention must be paid to such economic disparities, including attention to the distribution of high-quality teachers, if education is to improve for all students in the United States.}
  }
  {
    \VI{Tầm nhìn về toán học trung học phổ thông được trình bày trong ấn phẩm này phụ thuộc vào giáo viên có kiến thức vững vàng về nội dung toán học. Nếu thiếu nền tảng kiến thức này và sự tự tin mà nó mang lại, giáo viên sẽ gặp khó khăn trong việc đưa luận lên vị trí trung tâm của chương trình học. Thực tế cho thấy, nhiều tiền hơn, nhiều nguồn lực hơn và giáo viên tốt hơn thường được dành cho học sinh đến từ gia đình giàu có nhất, trong khi học sinh đến từ gia đình nghèo nhất lại chỉ nhận được ít tiền hơn, ít nguồn lực hơn và giáo viên kém chuẩn hơn (Darling-Hammond 2004). Chỉ riêng việc cung cấp nguồn lực công bằng cho tất cả học sinh chắc chắn chưa đủ để bảo đảm rằng mỗi học sinh sẽ phát triển được luận và hiểu trong toán học; tuy nhiên, đó chắc chắn là điều kiện cần để có thể đạt được tiến bộ thực chất (Bishop and Forgasz 2006). Vì vậy, nếu muốn cải thiện giáo dục cho tất cả học sinh tại Hoa Kì, cần phải quan tâm nghiêm túc đến chênh lệch kinh tế này, bao gồm cả việc chú ý đến sự phân bổ giáo viên chất lượng cao.}
  }
  {
    \NOTE{Đoạn này chốt lại lập luận về bất công kinh tế: nguồn lực công bằng không tự nó bảo đảm học sinh phát triển luận và hiểu, nhưng là điều kiện cần để tiến bộ thực chất có thể xảy ra. Trọng tâm không chỉ là tiền và cơ sở vật chất, mà còn là sự phân bổ giáo viên có chuyên môn vững vàng.}
  }

\TriPara
  {
    \EN{A positive trend is that more African American, Hispanic, and Native American students are studying more advanced mathematics than in previous decades. However, persistent gaps in course taking remain between these students and their white and Asian/Pacific Islander peers. In 2004 approximately 6 percent each of African American, Hispanic, and Native American high school students graduated having completed calculus, as compared with 33 percent of Asian/Pacific Islanders and 16 percent of white students (Planty, Provasnik, and Daniel 2007). The notion that African American, Hispanic, and Native American students are not as interested in or do not value taking advanced mathematics courses as much as their white peers is simply not true. Studies have shown that especially African American and Hispanic students “sometimes have more positive attitudes towards mathematics and higher educational aspirations” (Walker 2007, p. 48) than do their white peers, particularly in the early years of high school. Yet the disparities still exist, often tied to discrepancies in resources, which can have important consequences for students' achievement in mathematics. Hoffer, Rasinski, and Moore (1995), as cited in Tate and Rousseau (2002), found that when African American students and white students completed the same mathematics courses, they had comparable achievement gains. Tate and Rousseau go on to state, “These findings suggest that [many] of racial and SES differences in mathematics achievement in grades 9 through 12 are a product of the quality and the number of mathematics courses that African American, White, Hispanic, high- and low-SES students complete during high secondary school” (p. 276). This powerful finding makes clear the need for high schools to be diligent in supporting all students, especially those students who are often underrepresented in high-level mathematics classes, to study more advanced mathematics.}
  }
  {
    \VI{Xu hướng tích cực là ngày càng có nhiều học sinh Mĩ gốc Phi, học sinh Hispanic và học sinh người Mĩ bản địa theo học các khóa toán nâng cao hơn so với các thập niên trước. Tuy nhiên, khoảng cách dai dẳng trong việc theo học khóa toán vẫn tồn tại giữa nhóm học sinh này với bạn học da trắng và gốc Á/đảo Thái Bình Dương. Năm 2004, chỉ khoảng 6 phần trăm học sinh trung học phổ thông người Mĩ gốc Phi, học sinh Hispanic và học sinh người Mĩ bản địa tốt nghiệp sau khi đã hoàn thành giải tích, so với 33 phần trăm học sinh gốc Á/đảo Thái Bình Dương và 16 phần trăm học sinh da trắng (Planty, Provasnik, and Daniel 2007). Quan niệm cho rằng học sinh Mĩ gốc Phi, học sinh Hispanic và học sinh người Mĩ bản địa kém quan tâm hoặc không coi trọng khóa toán nâng cao như bạn da trắng là hoàn toàn không đúng. Nghiên cứu cho thấy, đặc biệt là học sinh Mĩ gốc Phi và học sinh Hispanic ``đôi khi có thái độ tích cực hơn đối với toán học và khát vọng giáo dục cao hơn'' (Walker 2007, p. 48) so với bạn da trắng, đặc biệt trong những năm đầu trung học phổ thông. Tuy vậy, chênh lệch vẫn tồn tại, thường gắn với chênh lệch về nguồn lực, và chênh lệch này có thể kéo theo hệ quả quan trọng đối với thành tích toán học của học sinh. Hoffer, Rasinski, and Moore (1995), được trích dẫn trong Tate and Rousseau (2002), phát hiện rằng khi học sinh Mĩ gốc Phi và học sinh da trắng hoàn thành cùng chương trình toán học, họ đạt được mức tăng trưởng thành tích tương đương. Tate and Rousseau tiếp tục nhận định: ``Phát hiện này cho thấy [phần lớn] khác biệt về chủng tộc và tình trạng kinh tế - xã hội (SES) trong thành tích toán học ở lớp 9 đến lớp 12 là hệ quả của chất lượng và số lượng khóa toán mà học sinh Mĩ gốc Phi, học sinh da trắng, học sinh Hispanic, học sinh thuộc SES cao và thấp hoàn thành trong giai đoạn trung học phổ thông'' (p. 276). Phát hiện có sức nặng này làm rõ nhu cầu các trường trung học phổ thông phải chủ động và nghiêm túc hỗ trợ tất cả học sinh, đặc biệt là những học sinh thường ít được đại diện trong các lớp toán trình độ cao, theo học các khóa toán nâng cao hơn.}
  }
  {
    \NOTE{ Đoạn này bác bỏ cách giải thích thiên lệch rằng học sinh Mĩ gốc Phi, học sinh Hispanic và học sinh người Mĩ bản địa ít quan tâm đến toán nâng cao. Lập luận chính là khi học sinh hoàn thành cùng khóa toán, mức tăng trưởng thành tích có thể tương đương; vì vậy, chênh lệch thành tích gắn chặt với cơ hội tiếp cận khóa toán chất lượng và số lượng khóa toán mà học sinh được học trong trung học phổ thông.}
  }

\TriPara
  {
    \EN{For more students from all racial and ethnic groups and socioeconomic levels to study more advanced mathematics, they must have the opportunity to enroll in these courses in their high school career. According to the Educational Testing Service (1999), schools having a majority of low-SES students offer fewer advanced mathematics courses than wealthier schools. This outcome, too, is closely connected with a lack of resources-most important, teachers who are qualified to teach advanced mathematics. Failure to offer advanced mathematics courses not only does a great disservice to the students but also perpetuates the shortage of people prepared to enter fields in science, technology, engineering, and mathematics.}
  }
  {
    \VI{Để có thêm nhiều học sinh thuộc mọi nhóm chủng tộc, sắc tộc và mọi mức kinh tế - xã hội theo học các khóa toán nâng cao, các em phải có cơ hội ghi danh vào những khóa học này trong quá trình học trung học phổ thông. Theo Educational Testing Service (1999), các trường có đa số học sinh thuộc nhóm ses thấp thường cung cấp ít khóa toán nâng cao hơn so với các trường có điều kiện hơn. Thực trạng này, một lần nữa, có mối liên hệ chặt chẽ với sự thiếu hụt nguồn lực, đặc biệt là giáo viên đủ chuẩn để dạy toán nâng cao. Việc không mở các khóa toán nâng cao không chỉ gây thiệt thòi lớn cho học sinh mà còn góp phần kéo dài tình trạng thiếu hụt người được chuẩn bị để bước vào các lĩnh vực khoa học, công nghệ, kĩ thuật và toán học.}
  }
  {
    \NOTE{Đoạn này nối cơ hội học toán nâng cao với công bằng dài hạn. Khi trường ít nguồn lực không mở được khóa toán nâng cao, thiệt thòi không chỉ dừng ở từng học sinh, mà còn kéo dài tình trạng thiếu người được chuẩn bị để bước vào các lĩnh vực khoa học, công nghệ, kĩ thuật và toán học.}
  }

\TriPara
  {
    \EN{Encouraging all students to take advanced mathematics courses is crucial. Walker (2007) offers suggestions that include (1) providing enrichment opportunities so that all students, but especially students from groups less represented in the mathematics pipeline, see mathematics as relevant to future careers and their lives; (2) involving families and communities; and (3) in schools that are less diverse, providing opportunities for engaging in mathematics in diverse student groups. As always, teachers play a crucial role in ensuring that once students are enrolled in advanced courses, those students are able to succeed. Teacher educators and staff developers must assist teachers in learning how to support mathematical reasoning in students who require additional support. For example, teachers can support the development of mathematical reasoning in English Language Learners (ELLs) by encouraging them to discuss their mathematical thinking on a daily basis and by providing them with various opportunities to communicate mathematically. Mathematical activity in which students are asked to reason is a vehicle for enhancing the linguistic skills of students who lack English vocabulary, thus avoiding their placement in remedial mathematics classes and preventing the delay of reasoning activities in the mathematics classroom until they have mastered basic linguistic skills (Moschkovich 2007). Instruction that centers on promoting group interaction, integrates language, uses multiple  representations, and emphasizes context is important for the success of ELLs (Bay-Williams and Herrera 2007).}
  }
  {
    \VI{Việc khuyến khích tất cả học sinh theo học các khóa toán nâng cao là điều hết sức quan trọng. Walker (2007) đưa ra những gợi ý, bao gồm: (1) cung cấp cơ hội bồi dưỡng mở rộng để tất cả học sinh, đặc biệt là học sinh thuộc những nhóm ít được đại diện trong lộ trình học toán, thấy toán học có liên quan đến nghề nghiệp tương lai và cuộc sống của các em; (2) huy động sự tham gia của gia đình và cộng đồng; và (3) ở những trường ít đa dạng hơn, tạo cơ hội cho học sinh tham gia hoạt động toán học trong các nhóm học sinh đa dạng. Giáo viên vẫn giữ vai trò then chốt trong việc bảo đảm rằng một khi học sinh đã ghi danh vào khóa học nâng cao, các em có thể học tập thành công. Người đào tạo giáo viên và người phụ trách phát triển chuyên môn cần hỗ trợ giáo viên học cách nâng đỡ luận toán học ở những học sinh cần hỗ trợ bổ sung. Chẳng hạn, giáo viên có thể hỗ trợ sự phát triển luận toán học ở học sinh đang học tiếng Anh (English Language Learners - ELLs) bằng cách khuyến khích các em thảo luận hằng ngày về tư duy toán học của mình và tạo cho các em nhiều cơ hội giao tiếp toán học. Hoạt động toán học trong đó học sinh được yêu cầu luận là phương tiện để tăng cường kĩ năng ngôn ngữ cho học sinh thiếu vốn từ tiếng Anh, qua đó tránh việc các em bị xếp vào lớp toán phụ đạo và tránh trì hoãn hoạt động luận trong lớp toán cho đến khi các em đã thành thạo kĩ năng ngôn ngữ cơ bản (Moschkovich 2007). Việc dạy học tập trung thúc đẩy tương tác nhóm, tích hợp ngôn ngữ, sử dụng nhiều biểu diễn và nhấn mạnh ngữ cảnh là quan trọng đối với sự thành công của học sinh ELL (Bay-Williams and Herrera 2007).}
  }
  {
    \NOTE{Đoạn này nhấn mạnh rằng khuyến khích học sinh vào khóa toán nâng cao phải đi cùng hỗ trợ để các em thành công sau khi đã ghi danh. Với học sinh học tiếng Anh, hoạt động toán học có luận không nên bị trì hoãn cho đến khi các em thành thạo tiếng Anh; ngược lại, chính việc thảo luận, giao tiếp và dùng nhiều biểu diễn trong toán học có thể hỗ trợ đồng thời cả luận toán học lẫn kĩ năng ngôn ngữ.}
  }

\ParallelSection
  {High Expectations}
  {Kì vọng cao}
  {}

\TriPara
  {
    \EN{In \emph{Principles and Standards for School Mathematics} (NCTM 2000a), the Equity Principle describes the importance of teachers' holding high expectations for all students. Teachers can motivate their students to perform at high levels by having and communicating high expectations, just as they can lower students' confidence in their mathematical abilities by exhibiting low expectations. Teachers' self-awareness of personal bias about who can and cannot do mathematics is an essential part of holding high expectations for all students. All too often, teachers' and administrators' views of students' ability, motivation, behavior, and future aspirations are influenced by their beliefs associated with such student identifiers as race, gender, socioeconomic status, native language, and home life. In turn, those beliefs can have serious consequences for the opportunities that a school provides to its students. By working to become aware of subconscious attitudes, teachers and administrators can strive to overcome those notions and act purposefully against social inequalities.}
  }
  {
    \VI{Trong \emph{Principles and Standards for School Mathematics} (NCTM 2000a), Nguyên lí Công bằng nhấn mạnh tầm quan trọng của việc giáo viên duy trì kì vọng cao đối với tất cả học sinh. Giáo viên có thể thúc đẩy học sinh đạt thành tích cao bằng cách duy trì và truyền đạt kì vọng cao, cũng giống như họ có thể làm suy giảm sự tự tin của học sinh vào năng lực toán học của bản thân nếu thể hiện kì vọng thấp. Việc giáo viên tự nhận thức về thiên lệch cá nhân liên quan đến việc ai có thể và ai không thể làm toán là phần thiết yếu của việc duy trì kì vọng cao đối với mọi học sinh. Quá thường xuyên, quan điểm của giáo viên và nhà quản lí về năng lực, động cơ, hành vi và khát vọng tương lai của học sinh bị chi phối bởi niềm tin gắn với những dấu hiệu nhận diện học sinh như chủng tộc, giới tính, tình trạng kinh tế - xã hội, ngôn ngữ mẹ đẻ và hoàn cảnh gia đình. Đến lượt mình, những niềm tin đó có thể gây ra hệ quả nghiêm trọng đối với cơ hội học tập mà nhà trường cung cấp cho học sinh. Bằng cách nỗ lực nhận diện những thái độ trong tiềm thức, giáo viên và nhà quản lí có thể phấn đấu vượt qua quan niệm đó và hành động có chủ đích nhằm chống lại bất bình đẳng xã hội.}
  }
  {
    \NOTE{Đoạn này chuyển từ công bằng về khóa học và nguồn lực sang công bằng trong kì vọng. Điểm chính là kì vọng của giáo viên không trung tính: chúng có thể nâng đỡ hoặc làm suy giảm niềm tin học toán của học sinh, nhất là khi bị chi phối bởi thiên lệch gắn với chủng tộc, giới tính, tình trạng kinh tế - xã hội, ngôn ngữ mẹ đẻ hoặc hoàn cảnh gia đình.}
  }

\TriPara
  {
    \EN{Although holding high expectations for all students is crucial, expectations alone are not enough. By taking specific action in the classroom, teachers can truly support students in meeting those high expectations. One such action is to encourage all students to share their thinking and listen to the thinking of others. Establishing a community of learning in which a multitude of approaches and solutions are encouraged and respected is a way to effect positive change (NCTM 2008). Using tasks with multiple entry points is one way to encourage this student interaction, as all students can make a meaningful contribution to the class discussion. Subtle teacher behaviors also demonstrate the belief that all students can be successful. For example, Cousins-Cooper (2000) suggested that teachers can encourage African American students when they get stuck on a problem by asking probing questions instead of showing the students how to do the problem; such behavior is important for teachers of all students. This behavior communicates to the students and all who observe the interaction that the teacher believes that her or his students can be successful in making sense of the problem and working toward a solution. Teachers must be constantly aware of the messages they communicate to students both verbally and nonverbally.}
  }
  {
    \VI{Mặc dù việc giữ kì vọng cao đối với tất cả học sinh là hết sức quan trọng, chỉ riêng kì vọng thôi thì chưa đủ. Thông qua hành động cụ thể trong lớp học, giáo viên mới có thể thực sự hỗ trợ học sinh đáp ứng những kì vọng cao đó. Một hành động như vậy là khuyến khích mọi học sinh chia sẻ cách nghĩ của mình và lắng nghe cách nghĩ của người khác. Việc xây dựng cộng đồng học tập trong đó nhiều cách tiếp cận và lời giải khác nhau được khuyến khích và tôn trọng là một cách đem lại thay đổi tích cực (NCTM 2008). Sử dụng nhiệm vụ có nhiều điểm vào là một cách để khuyến khích tương tác này giữa học sinh, bởi mọi học sinh đều có thể đóng góp một cách có ý nghĩa cho thảo luận chung của lớp. Những hành vi tinh tế của giáo viên cũng thể hiện niềm tin rằng tất cả học sinh đều có thể thành công. Chẳng hạn, Cousins-Cooper (2000) gợi ý rằng giáo viên có thể khích lệ học sinh Mĩ gốc Phi khi các em gặp bế tắc với một bài toán bằng cách đặt câu hỏi thăm dò thay vì chỉ cho các em cách làm; cách ứng xử này cũng quan trọng đối với giáo viên dạy mọi học sinh. Cách ứng xử như vậy truyền đạt tới học sinh, và tới tất cả người quan sát tương tác đó, rằng giáo viên tin tưởng học sinh của mình có khả năng hiểu bài toán và tiến tới lời giải. Do đó, giáo viên cần liên tục ý thức về thông điệp mà họ truyền đạt tới học sinh, cả bằng lời nói lẫn bằng tín hiệu phi ngôn ngữ.}
  }
  {
    \NOTE{Đoạn này nhấn mạnh rằng kì vọng cao phải được chuyển thành hành động cụ thể trong lớp học. Những hành động nhỏ như yêu cầu học sinh chia sẻ cách nghĩ, lắng nghe cách nghĩ của bạn, dùng nhiệm vụ có nhiều điểm vào, hoặc đặt câu hỏi thăm dò khi học sinh gặp bế tắc đều truyền đi thông điệp rằng mọi học sinh đều có thể tham gia có ý nghĩa vào toán học.}
  }

\TriPara
  {
    \EN{Creating an inclusive environment in which all students-regardless of race, gender, socioeconomic status, or perceived motivation and learning ability-are engaged in reasoning and sense-making activities communicates to students a belief that everyone is expected to be successful in mathematics. Researchers have noted a tremendous difference in success rates among highminority, low-socioeconomic student populations when a culture of high expectations and a shared goal of success for all students is adopted not just by individual teachers but by the entire school, including families and the community (Gutierrez 2000; Education Trust 2005).}
  }
  {
    \VI{Việc tạo dựng môi trường hòa nhập, trong đó mọi học sinh, bất kể chủng tộc, giới tính, tình trạng kinh tế - xã hội, hay động cơ và khả năng học tập được người khác nhìn nhận như thế nào, đều tham gia vào hoạt động luận và hiểu sẽ truyền đạt tới học sinh niềm tin rằng mọi học sinh đều được kì vọng sẽ thành công trong toán học. Các nhà nghiên cứu đã ghi nhận những khác biệt rất lớn về tỉ lệ thành công ở các nhóm học sinh có tỉ lệ thiểu số cao và điều kiện kinh tế - xã hội thấp, khi văn hóa kì vọng cao và mục tiêu chung về thành công của mọi học sinh được tiếp nhận không chỉ bởi từng giáo viên riêng lẻ mà bởi toàn bộ nhà trường, bao gồm cả gia đình và cộng đồng (Gutierrez 2000; Education Trust 2005).}
  }
  {
    \NOTE{Đoạn này mở rộng kì vọng cao từ hành vi của từng giáo viên sang văn hóa chung của cả nhà trường. Một môi trường hòa nhập chỉ thật sự có tác dụng khi mọi học sinh được tham gia vào hoạt động luận và hiểu, đồng thời gia đình, cộng đồng và toàn bộ nhà trường cùng chia sẻ mục tiêu rằng mọi học sinh đều có thể thành công trong toán học.}
  }

\ParallelSection
  {Conclusion}
  {Kết luận}
  {}

\TriPara
  {
    \EN{In general, teachers and schools across North America are working tirelessly toward the goal of providing high-quality education to every student in their classrooms. However, in the shuffle of high school life, some students are inadvertently overlooked or underserved. Moreover, the same students are frequently among those overlooked-those of color, those of low or exceptionally high achievement, and those of low socioeconomic status. Teachers, teacher leaders, and administrators, along with families and the community, must constantly reflect on how educational policy and practices may work to support or fail to support giving all students the opportunities with reasoning and sense making that they need to be successful as they exit high school. Some questions to reflect on in trying to reach this goal include these: Are the students enrolled in advanced mathematics courses at a school representative of the demographics of the school? If a school employs tracking, do students move between tracks flexibly and with the support of teachers and administrators? Are decisions on course placements based on evidence rather than subjective judgments that may unintentionally reinforce group stereotypes? In any given mathematics course, are students engaged in reasoning and sense making on a daily basis?}
  }
  {
    \VI{Nhìn chung, giáo viên và nhà trường trên khắp Bắc Mĩ đang làm việc không mệt mỏi nhằm hướng tới mục tiêu cung cấp nền giáo dục chất lượng cao cho mọi học sinh trong lớp học. Tuy nhiên, giữa guồng quay bận rộn của đời sống trung học phổ thông, một số học sinh vẫn vô tình bị bỏ sót hoặc không được phục vụ đầy đủ. Hơn nữa, những học sinh bị bỏ sót thường vẫn là cùng những nhóm ấy: học sinh da màu, học sinh có thành tích thấp hoặc đặc biệt cao, và học sinh thuộc nhóm kinh tế - xã hội thấp. Giáo viên, giáo viên giữ vai trò dẫn dắt chuyên môn, nhà quản lí, cùng với gia đình và cộng đồng, cần liên tục suy ngẫm về cách chính sách và thực hành giáo dục có thể hỗ trợ, hoặc không hỗ trợ, việc bảo đảm cho tất cả học sinh có được cơ hội luận và hiểu cần thiết để thành công khi hoàn thành bậc trung học phổ thông. Một số câu hỏi gợi suy ngẫm để hướng tới mục tiêu này bao gồm: Học sinh ghi danh vào các khóa toán nâng cao của nhà trường có đại diện cho cơ cấu nhân khẩu học của nhà trường hay không? Nếu nhà trường áp dụng phân luồng, học sinh có được di chuyển giữa các luồng một cách linh hoạt và có sự hỗ trợ của giáo viên và nhà quản lí hay không? Quyết định xếp học sinh vào khóa học có dựa trên bằng chứng, thay vì phán đoán chủ quan có thể vô tình củng cố khuôn mẫu nhóm, hay không? Trong bất kì khóa toán cụ thể nào, học sinh có được tham gia vào hoạt động luận và hiểu hằng ngày hay không?}
  }
  {
    \NOTE{Phần kết luận gom lại ba trục công bằng đã triển khai trong chương: ai được vào khóa toán nâng cao, ai bị giữ trong luồng thấp, quyết định xếp lớp dựa trên bằng chứng hay thiên lệch, và liệu mọi học sinh có thật sự được tham gia vào luận và hiểu hằng ngày hay không. Trọng tâm là biến công bằng thành câu hỏi tự kiểm tra thường trực của giáo viên, nhà trường, gia đình và cộng đồng.}
  }

\CloseGrid